Sei $\ell$ die Länge und $m$ die Masse des Bretts, $b$ der Abstand der Stütze vom Bolzen und $M$ die Masse des Springers. Die Gewichtskraft des Bretts greift in seiner Mitte an.

Nimm den Bolzen als Drehachse: Die Stütze dreht das Brett nach oben, die Gewichtskräfte des Bretts und des Springers drehen es nach unten:
\begin{align}
 F_\mathrm{S}\cdot b &= m\,g\cdot\frac{\ell}{2} + M\,g\cdot\ell
\end{align}
Also
\begin{align}
 F_\mathrm{S} &= \FStF \\
   &= \frac{\mB\cdot\g\cdot\frac{\LB}{2}+\mS\cdot\g\cdot\LB}{\bS} \\
   &= \FSt \approx \result{\FStP{3}{2}}
\end{align}
Dann die Kräfte: Nach oben wirkt die Stütze, nach unten der Bolzen und die beiden Gewichtskräfte.
\begin{align}
 F_\mathrm{B} &= \FBzF \\
   &= \FSt-(\mB+\mS)\cdot\g \\
   &= \FBz \approx \result{\FBzP{3}{2}}
\end{align}
Die Stütze trägt ein Vielfaches der Gewichtskraft des Springers: Das Brett wirkt als Hebel.
